%continuation of last week's proof
\todo{Lecture 25}
Soient $\{ w_{1}^{i},\dots,w_{\ell_{i}}^{i} \}$ une base de $W_{i}$. Alors $P=\{ w_{1}^{1},\dots,w_{\ell_{1}}^{1},\dots,w_{1}^{k},\dots,w_{\ell _{k}} \}\in \mathbf{C}^{n\times n}$ et inversible et 
$$
P^{-1}AP=\mathrm{diag}(B_{1},\dots ,B_{k})=D
$$
est une matrice bloc diagonale ou $(B_{i}-\lambda_{i}I_{\ell _{i}})$ est une matrice nilpotente pour chaque $B_{i}\in \mathbf{C}^{\ell _{i}\times \ell _{i}}$.
\begin{question}
Est-ce que $\ell _{i}=n_{i}$ pour tout $i\in[k]$ ?
\end{question}
\begin{answer}
On a 
\begin{align*}
Je_{i} & =[Aw_{\mu}^{j}]_{W} \\
   & =[(A-\lambda _{i}I)w_{\mu}^{j}]_{W} \\
 & =(D-\lambda _{i}I)e_{i}\\
& =\begin{pmatrix}
\boldsymbol{0} \\
\mathrm{col}(B_{j})_{\mu} \\
\boldsymbol{0}
\end{pmatrix}.
\end{align*}
On a
$$
P^{-1}(A-\lambda _{i}I)P=P^{-1}AP-\lambda _{i}I.
$$
\end{answer}
\end{proof}
\begin{example}
Soit
$$
A=\begin{pmatrix}
25 & 34 & 18 \\
-14 & -19 & -10 \\
-4 & -6 & -1 
\end{pmatrix}
$$
de polynôme caractéristique
$$
f_{A}(x)=(x-1)^{2}(x-3).
$$
Alors
$$
\lambda_{1}=1\Rightarrow n_{1}=2,\quad \lambda_{2}=3\Rightarrow n_{2}=1.
$$
On calcule
$$
\ker(A-1I)^{2}=\mathrm{span}\{ \begin{pmatrix}
2 \\
0 \\
-1
\end{pmatrix}, \begin{pmatrix}
0 \\
2 \\
-1
\end{pmatrix} \}
$$
et
$$
\ker(A-3I)^{1}=\mathrm{span}\{ \begin{pmatrix}
-7 \\
4 \\
1
\end{pmatrix}\}.
$$
On construit $P$ :
$$
P=\begin{pmatrix}
2 & 0 & -7 \\
0 & 2 & 4 \\
-1 & -1  & 1
\end{pmatrix}
$$
et
$$
P^{-1}AP=\begin{pmatrix}
16 & 25 & 0 \\
-9 & -14 & 0 \\
0 & 0 & 3
\end{pmatrix}.
$$
On vérifie que les $(B_{i}-\lambda _{i}I_{n_{i}})^{n_{i}}$ sont nilpotentes
$$
\underbrace{ \begin{pmatrix}
15 & 25 \\
-9 & -15
\end{pmatrix}^{2} }_{ N_{1} }=0, \quad N_{2}=0=0.
$$
Donc
$$
P^{-1}AP=\begin{pmatrix}
N_{1}+1I_{2} \\
 & N_{2}+3I_{1}
\end{pmatrix}
$$
ou $N_{1}\in \mathbf{C}^{2\times2}$ et $N_{2}\in \mathbf{C}^{1\times 1}$ sont nilpotentes.
\end{example}
\begin{reminder}
$$
A\in \mathbf{C}^{n\times n}\sim \mathrm{diag}(N_{1}+\lambda_{1}I_{1},\dots,N_{k}+\lambda _{k}I_{k}).
$$
\end{reminder}

On veut montrer que si $N\in K^{n\times n}$ est nilpotente, alors il existe $P\in K^{n\times n}$ inversible telle que $P^{-1}NP=J$, ou $J$ est une forme normale de Jordan avec tous les elements sur la diagonale nuls. Ceci est une consequence du théorème suivant.
\begin{theorem}
Soit $V$ un espace vectoriel sur $K$ et $N:V\to V$ un endomorphisme nilpotente. Alors, ils existent $v_{1},\dots,v_{k}\in V$ tels que :
\begin{enumerate}
\item $v_{1},N(v_{1}),\dots,N^{\ell_{1}}(v_{1}),\dots,v_{k},N(v_{k}),\dots,N^{\ell _{k}}(v_{k})$ forment une base de $V$. 
\item Pour tous $j\in[k]$, on a que $N^{\ell _{j}+1}(v_{j})=0$.
\end{enumerate}
\end{theorem}
\begin{remark}
Soit $B'=\{ N^{\ell _{k}}(v_{k}),\dots,v_{k},\dots,N^{\ell_{1}}(v_{1}),\dots,v_{1} \}$. Alors
$$
A^{N}_{B'}=\begin{pmatrix}
J(0,\ell _{k}) \\
 & J(0,\ell _{k-1}) \\
 &  & \ddots \\
 &  &  & J(0,\ell_{1})
\end{pmatrix}\in \mathbf{C}^{n\times n}. 
$$
\end{remark}
\begin{proof}
Soit $i\in \mathbf{N}$ maximal tel que $N^{i}\neq 0$. Si $i=1$. Si $i=1$ $\checkmark$. Sinon, il existe $v\in V$ tel que $N^{i}(v)\neq0$. On veut montrer qu'on peut completer les vecteurs $v,N(v),\dots,N^{i}(v)$ avec $w_{1},\dots,w_{n-(i+1)}$ pour former une base
$$
B=\{ v,N(v),\dots,N^{i}(v),w_{1},\dots,w_{n-(i+1)} \},
$$
telle que $N(w_{i})\in \mathrm{span}\{ w_{1},\dots,w_{n-(i+1)} \}=W$.

Par recurrence, le théorème est correct pour $N:W\to W$. Les vecteurs $v,N(v),\dots,N^{i}(v)$ sont linéairement independents. Imaginons que ce ne soit pas le cas, alors
$$
0v+0N(v)+3N^{2}(v)+\dots=0,
$$
qui implique que $N^{2}(v)$ est une combinaison linéaire de $N^{3}(v),..,N^{i}(v)$ mais comme $N^{i}(v)\neq 0$ on a que 
$$
N^{i-2}(N^{2}(v))\neq 0,
$$
qui est une contradiction. 

On va supposer, qu'ils existent $x,y\in V$ tels que 
$$
B=\{v,N(v),\dots,N^{i}(v),x,y\}
$$
soit une base de $V$. On peut choisir ces vecteurs tels que 
$$
[N(x)]_{B}=\begin{pmatrix}
\alpha^{x} \\
\boldsymbol{0} \\
\beta^{x} \\
\gamma^{x}
\end{pmatrix},\quad [N(y)]_{B}=\begin{pmatrix}
\alpha^{y} \\
\boldsymbol{0} \\
\beta^{y} \\
\gamma^{y}
\end{pmatrix}.
$$
Si $\alpha\neq 0$, alors
$$
N(x)=\alpha^{x}v+\beta^{x}x+\gamma^{x}y,
$$
donc
$$
N^{2}(x)=\delta v+\alpha^{x}N(x)+\underbrace{ \eta }_{ \in \mathrm{span}\{ x,y \} }
$$
et
$$
N^{i+1}(x)=\underbrace{ \eta }_{ \in \mathrm{span}\{ v,\dots,N^{i-1}(v) \} } +\alpha^{x}N^{i}(v)+\underbrace{ \mu }_{ \in \mathrm{span}\{ x,y \} }\neq 0.
$$
Mais $N^{i+1}(x)=0$ donc on a une contradiction. 

\end{proof}